\section*{Abstract}

Decentralized Finance (DeFi) has experienced exponential growth since 2020, with lending protocols such as Aave, Compound, and MakerDAO intermediating total value locked exceeding \$100 billion across multiple blockchain networks. These protocols operate entirely through smart contracts that govern collateralization, borrowing, and liquidation automatically, creating a fully transparent on-chain record of all participant behavior. However, the academic literature on DeFi lending has predominantly focused on protocol-level risk exposure, liquidation mechanics, and aggregate market efficiency, devoting considerably less attention to the micro-level behavioral responses of individual borrowers as their positions approach risky thresholds. This research report surveys the relevant literature across three interconnected domains—DeFi lending infrastructure, behavioral finance and prospect theory, and on-chain credit risk assessment—and identifies a specific, underexplored research gap at their intersection. We propose a study that investigates whether borrowers' active adjustments during periods of rising position risk provide risk information beyond that captured by conventional on-chain indicators such as health factor, collateralization ratio, and account activity. The proposed methodology reconstructs position-risk trajectories from publicly available blockchain data, classifies borrower-initiated actions separately from passive liquidation events, and tests the incremental predictive contribution of behavioral process variables using a discrimination framework. A discussion of possible outcomes, including contributions to DeFi protocol design and credit risk modeling, as well as a critical assessment of the study's boundaries, is provided.

\vspace{8pt}
\noindent\textbf{Keywords:} Decentralized Finance, DeFi lending, behavioral finance, prospect theory, credit risk, blockchain analytics, on-chain data
